\chapter{วิธีการดำเนินการวิจัย}
\label{chapter:experiment}

บทนี้อธิบายระเบียบวิธีที่ใช้ในการดำเนินโครงงาน ประกอบด้วยภาพรวมของกระบวนการทำงาน การเตรียมชุดข้อมูล การแบ่งชุดข้อมูลฝึกและชุดทดสอบ การพัฒนาแบบจำลอง เกณฑ์การประเมินผล การอธิบายผลของแบบจำลอง การพัฒนาระบบต้นแบบ และมาตรการที่ใช้ควบคุมให้การทดลองสามารถทำซ้ำได้

\section{ภาพรวมของกระบวนการดำเนินงาน}

การดำเนินงานแบ่งออกเป็น 5 ระยะที่เชื่อมต่อกันเป็นไปป์ไลน์เดียว โดยผลลัพธ์ของแต่ละระยะเป็นอินพุตของระยะถัดไป ดังแสดงในรูปที่ \ref{fig:workflow}

\begin{figure}[ht]
    \centering
    \begin{tikzpicture}[
        node distance = 7mm,
        box/.style = {
            draw, rounded corners=2pt, thick,
            text width=88mm, align=center,
            minimum height=10mm, inner sep=2mm
        },
        arr/.style = {-latex, thick}
    ]
        \node[box] (p1) {ระยะที่ 1: การศึกษาและเตรียมข้อมูล};
        \node[box, below=of p1] (p2) {ระยะที่ 2: การพัฒนาแบบจำลอง};
        \node[box, below=of p2] (p3) {ระยะที่ 3: การประเมินผลและเปรียบเทียบ};
        \node[box, below=of p3] (p4) {ระยะที่ 4: การอธิบายผลของแบบจำลอง};
        \node[box, below=of p4] (p5) {ระยะที่ 5: การแสดงผลและสรุปผล};
        \draw[arr] (p1) -- (p2);
        \draw[arr] (p2) -- (p3);
        \draw[arr] (p3) -- (p4);
        \draw[arr] (p4) -- (p5);
    \end{tikzpicture}
    \caption{ภาพรวมกระบวนการดำเนินงานของโครงงาน}
    \label{fig:workflow}
\end{figure}

\section{การศึกษาและเตรียมข้อมูล}

\subsection{ชุดข้อมูลที่ใช้}

โครงงานนี้ใช้ชุดข้อมูล CICIDS2017 รุ่น \textit{MachineLearningCVE} เป็นคอร์ปัสหลักเพียงชุดเดียว ชุดข้อมูลนี้ประกอบด้วยไฟล์บันทึกการจับข้อมูล (capture file) จำนวน 8 ไฟล์ ครอบคลุมการดำเนินการโจมตีตามตารางเวลาที่กำหนดไว้ล่วงหน้าเป็นเวลา 5 วันทำการ \cite{sharafaldin2018cicids2017} รายละเอียดของไฟล์และประเภทการโจมตีที่ปรากฏในแต่ละไฟล์แสดงในตารางที่ \ref{tab:capture-files}

\begin{table}[ht]
    \caption{ไฟล์บันทึกการจับข้อมูลในชุดข้อมูล CICIDS2017 และประเภทการโจมตีที่ปรากฏ}
    \label{tab:capture-files}
    \small
    \begin{tabularx}{\linewidth}{@{}l X l@{}}
        \hline
        \textbf{วันที่บันทึก} & \textbf{ชื่อไฟล์} & \textbf{ประเภทการโจมตี} \\
        \hline
        จันทร์ & Monday-WorkingHours & ไม่มี (BENIGN) \\
        อังคาร & Tuesday-WorkingHours & Brute Force \\
        พุธ & Wednesday-workingHours & DoS, Heartbleed \\
        พฤหัสบดี (เช้า) & Thursday-WorkingHours-Morning-WebAttacks & Web Attack \\
        พฤหัสบดี (บ่าย) & Thursday-WorkingHours-Afternoon-Infilteration & Infiltration \\
        ศุกร์ (เช้า) & Friday-WorkingHours-Morning & Bot \\
        ศุกร์ (บ่าย) & Friday-WorkingHours-Afternoon-PortScan & PortScan \\
        ศุกร์ (บ่าย) & Friday-WorkingHours-Afternoon-DDos & DDoS \\
        \hline
    \end{tabularx}
\end{table}

การใช้ไฟล์ทั้ง 8 ไฟล์เป็นเงื่อนไขบังคับ เนื่องจากการโจมตีแต่ละตระกูลปรากฏอยู่ในไฟล์เพียงไฟล์เดียว หากไฟล์ใดขาดหายไปจะทำให้คลาสนั้นหายไปจากชุดข้อมูลทั้งคลาส กระบวนการโหลดข้อมูลจึงได้รับการออกแบบให้ตรวจสอบความครบถ้วนของไฟล์และหยุดการทำงานพร้อมแจ้งข้อผิดพลาด แทนที่จะสร้างแคชข้อมูลที่ไม่สมบูรณ์

\subsection{การสำรวจข้อมูลเบื้องต้น}

ก่อนการเตรียมข้อมูล คณะผู้จัดทำได้ดำเนินการสำรวจข้อมูลเบื้องต้น (Exploratory Data Analysis) เพื่อทำความเข้าใจโครงสร้างของชุดข้อมูล ประกอบด้วยการตรวจสอบการกระจายตัวของคลาส การตรวจสอบค่าว่างและค่าอนันต์ในแต่ละคอลัมน์ การวิเคราะห์สหสัมพันธ์ระหว่างคุณลักษณะ และการพิจารณาการกระจายตัวของคุณลักษณะเชิงตัวเลข ผลการสำรวจยืนยันความไม่สมดุลอย่างรุนแรงของคลาส ซึ่งเป็นข้อมูลนำเข้าสำคัญต่อการออกแบบกลยุทธ์การจัดการความไม่สมดุลในระยะถัดไป

\subsection{การทำความสะอาดข้อมูล}

ขั้นตอนการทำความสะอาดข้อมูลประกอบด้วยการดำเนินการต่อไปนี้

\begin{enumerate}
    \item ตัดอักขระช่องว่างส่วนเกินออกจากชื่อคอลัมน์ทั้งหมด เนื่องจากไฟล์ต้นฉบับมีชื่อคอลัมน์ที่ขึ้นต้นด้วยช่องว่าง
    \item ลบคอลัมน์ Fwd Header Length.1 ซึ่งเป็นคอลัมน์ที่ซ้ำซ้อนกับคอลัมน์ Fwd Header Length และเป็นข้อบกพร่องที่ทราบแล้วของชุดข้อมูลรุ่นที่เผยแพร่
    \item ลบระเบียนที่มีค่าว่าง (missing value) และค่าอนันต์ (infinity) ในคุณลักษณะที่ใช้เป็นอินพุตของแบบจำลอง
    \item ลบระเบียนที่มีป้ายกำกับที่ไม่สามารถจับคู่เข้ากับคลาสมาตรฐานทั้ง 9 คลาสได้
    \item บันทึกตำแหน่งของแถวในไฟล์ต้นฉบับไว้ในคอลัมน์เมทาดาทา \_row\_index ก่อนการลบระเบียนใด ๆ เพื่อใช้เป็นกุญแจลำดับการจับข้อมูลในขั้นตอนการแบ่งข้อมูล
\end{enumerate}

หลังการทำความสะอาด ชุดข้อมูลที่ใช้งานประกอบด้วยระเบียนโฟลว์จำนวน 2,497,980 ระเบียน โดยมีคุณลักษณะเชิงตัวเลขที่ใช้เป็นอินพุตของแบบจำลองจำนวน 77 คุณลักษณะ ส่วนคอลัมน์เมทาดาทา ได้แก่ \text{dataset\_id}, \text{source\_file}, \text{capture\_window} และ \text{\_row\_index} ถูกกันไว้ไม่ให้เข้าสู่แบบจำลอง เพื่อป้องกันการรั่วไหลของสารสนเทศที่บ่งชี้คลาสโดยตรง

\subsection{การปรับป้ายกำกับให้เป็นมาตรฐาน}

ป้ายกำกับดิบในชุดข้อมูล CICIDS2017 มีทั้งสิ้น 15 ค่า ซึ่งบางค่ามีความละเอียดในระดับเครื่องมือที่ใช้โจมตีมากกว่าระดับตระกูลของการโจมตี และบางค่ามีอักขระควบคุมที่ไม่ใช่รหัสแอสกีปนอยู่ในสายอักขระ กระบวนการปรับป้ายกำกับให้เป็นมาตรฐานจึงดำเนินการตามลำดับดังนี้ ทำการนอร์มัลไลซ์ยูนิโคดแบบ NFKD แทนที่อักขระควบคุมด้วยช่องว่างหนึ่งตัว รวบช่องว่างที่ติดกันให้เหลือตัวเดียว แปลงเป็นตัวพิมพ์เล็ก และตัดช่องว่างหัวท้ายออก ผลลัพธ์ที่ได้ถูกนำไปจับคู่เข้ากับคลาสมาตรฐาน 9 คลาสตามตารางที่ \ref{tab:label-mapping}

\begin{table}[ht]
    \caption{การจับคู่ป้ายกำกับดิบในชุดข้อมูล CICIDS2017 เข้ากับคลาสมาตรฐาน 9 คลาส}
    \label{tab:label-mapping}
    \small
    \begin{tabularx}{\linewidth}{@{}X l@{}}
        \hline
        \textbf{ป้ายกำกับดิบ} & \textbf{คลาสมาตรฐาน} \\
        \hline
        BENIGN & BENIGN \\
        DoS Hulk, DoS GoldenEye, DoS slowloris, DoS Slowhttptest & DoS \\
        DDoS & DDoS \\
        PortScan & PortScan \\
        FTP-Patator, SSH-Patator & Brute Force \\
        Web Attack -- Brute Force, Web Attack -- XSS, Web Attack -- Sql Injection & Web Attack \\
        Bot & Bot \\
        Infiltration & Infiltration \\
        Heartbleed & Heartbleed \\
        \hline
    \end{tabularx}
\end{table}

เพื่อรองรับป้ายกำกับรูปแบบใหม่ที่อาจเพิ่มเข้ามาในอนาคต กระบวนการจับคู่ยังมีตารางสำรองที่พิจารณาจากคำนำหน้าของป้ายกำกับ เช่น ป้ายกำกับที่เริ่มต้นด้วย \texttt{dos} จะถูกจับคู่เข้ากับคลาส DoS โดยอัตโนมัติ จากนั้นป้ายกำกับที่เป็นข้อความจะถูกแปลงเป็นจำนวนเต็มด้วยการเข้ารหัสป้ายกำกับ (label encoding) เพื่อให้แบบจำลองนำไปประมวลผลได้

\subsection{การจัดการคุณลักษณะ}

การปรับสเกลของคุณลักษณะดำเนินการด้วยตัวปรับมาตรฐาน (standard scaler) ซึ่งแปลงคุณลักษณะแต่ละตัวให้มีค่าเฉลี่ยเป็นศูนย์และส่วนเบี่ยงเบนมาตรฐานเป็นหนึ่ง ตามสมการ

\begin{equation}
    z_{j} = \frac{x_{j} - \mu_{j}}{\sigma_{j}}
\end{equation}

โดยที่ $\mu_{j}$ และ $\sigma_{j}$ คือค่าเฉลี่ยและส่วนเบี่ยงเบนมาตรฐานของคุณลักษณะที่ $j$ ซึ่งประมาณค่าจากชุดฝึกเท่านั้น

ประเด็นสำคัญเชิงระเบียบวิธีคือ ตัวปรับสเกลถูกประกอบเข้าเป็นขั้นตอนแรกของไปป์ไลน์ร่วมกับตัวจำแนกประเภท มิได้ปรับสเกลข้อมูลไว้ล่วงหน้าแยกจากกระบวนการฝึก ด้วยวิธีนี้ ในการตรวจสอบไขว้และการค้นหาไฮเปอร์พารามิเตอร์ทุกครั้ง ค่าสถิติของตัวปรับสเกลจะถูกประมาณใหม่จากข้อมูลในโฟลด์ฝึกเท่านั้น จึงไม่เกิดการรั่วไหลของข้อมูลจากโฟลด์ตรวจสอบเข้าสู่กระบวนการฝึก \cite{pedregosa2011scikit}

สำหรับการคัดเลือกคุณลักษณะ โครงงานนี้เลือกใช้คุณลักษณะทั้ง 77 ตัวเป็นอินพุตของแบบจำลอง แทนการคัดกรองล่วงหน้า ด้วยเหตุผลสองประการ ประการแรก แบบจำลองที่อาศัยต้นไม้ตัดสินใจมีกลไกคัดเลือกคุณลักษณะอยู่ในตัวผ่านการเลือกจุดแบ่ง ประการที่สอง การคัดกรองคุณลักษณะโดยพิจารณาความสัมพันธ์กับป้ายกำกับจากทั้งชุดข้อมูลก่อนการแบ่งข้อมูลเป็นรูปแบบหนึ่งของการรั่วไหลของข้อมูลที่พบได้บ่อย \cite{arp2022dosdonts} การวิเคราะห์ความสำคัญของคุณลักษณะจึงกระทำภายหลังการฝึกด้วยเทคนิค SHAP ตามที่กล่าวถึงในหัวข้อ \ref{sec:shap}

\section{การแบ่งชุดข้อมูลฝึกและชุดทดสอบ}
\label{sec:split}

\subsection{หลักการแบ่งข้อมูลตามลำดับการจับข้อมูล}

โครงงานนี้ใช้การแบ่งข้อมูลตามลำดับการจับข้อมูล (capture-order split) ในสัดส่วน 70/30 ภายในทุกกลุ่มของคู่ (ไฟล์การจับข้อมูล, คลาส) กล่าวคือ ภายในแต่ละกลุ่ม ระเบียน 70\% แรกตามลำดับแถวในไฟล์ถูกจัดเป็นชุดฝึก และระเบียน 30\% ท้ายถูกจัดเป็นชุดทดสอบ เงื่อนไขนี้ทำให้ภายในกลุ่มเดียวกัน ไม่มีระเบียนในชุดทดสอบใดอยู่ก่อนระเบียนในชุดฝึกตามลำดับการจับข้อมูล และลดโอกาสที่ระเบียนซึ่งอยู่ติดกันและมีลักษณะใกล้เคียงกันจะกระจายไปอยู่ทั้งสองฝั่ง อย่างไรก็ตาม การแบ่งนี้ไม่ได้ตัดที่จุดเวลาเดียวกันทั้งคอร์ปัส จึงไม่ได้จำลองสถานการณ์ที่ฝึกด้วยข้อมูลในอดีตทั้งหมดแล้วทำนายข้อมูลในอนาคตอย่างสมบูรณ์

เนื่องจากชุดข้อมูลรุ่น \textit{MachineLearningCVE} ไม่มีคอลัมน์เวลา กุญแจลำดับที่ใช้จึงเป็นตำแหน่งของแถวในไฟล์ต้นฉบับ มิใช่เวลาประทับ (timestamp) ลำดับนี้ได้รับการตรวจสอบความสอดคล้องกับตารางเวลาการโจมตีที่ผู้จัดทำชุดข้อมูลเผยแพร่ไว้ทุกครั้งก่อนการฝึก ข้อควรระวังที่ต้องบันทึกไว้คือ เครื่องมือ CICFlowMeter สร้างระเบียนเมื่อโฟลว์สิ้นสุดลง ลำดับแถวจึงน่าจะสะท้อนลำดับที่โฟลว์สิ้นสุด ไม่ใช่ลำดับที่โฟลว์เริ่มต้น ซึ่งสอดคล้องกับระบบตรวจจับการบุกรุกที่ทำงานกับคุณลักษณะระดับโฟลว์ เพราะคุณลักษณะเหล่านั้นคำนวณได้เมื่อโฟลว์สิ้นสุดแล้วเท่านั้น

\subsection{ข้อจำกัดที่ทำให้ไม่สามารถกันข้อมูลทั้งไฟล์ได้}

การกันข้อมูลทั้งไฟล์ไว้เป็นชุดทดสอบ (source holdout) ซึ่งเป็นแนวทางที่เข้มงวดกว่า ไม่สามารถทำได้บนคอร์ปัสนี้ เพราะการโจมตีแต่ละคลาสปรากฏอยู่ในไฟล์เดียว การกันไฟล์หนึ่งออกไปจะทำให้คลาสในไฟล์นั้นไม่มีตัวอย่างในชุดฝึกเลย ด้วยเหตุผลเดียวกัน การตรวจสอบไขว้จึงใช้วิธี Stratified $k$-Fold แทน Group $k$-Fold ที่จัดกลุ่มตามไฟล์ต้นทาง ข้อจำกัดนี้ถูกบันทึกไว้อย่างชัดเจนเพื่อประกอบการตีความผลลัพธ์

\subsection{จำนวนตัวอย่างในแต่ละคลาส}

ผลของการแบ่งข้อมูลตามหลักการข้างต้นทำให้ได้ชุดฝึกจำนวน 1,748,579 ระเบียน และชุดทดสอบจำนวน 749,401 ระเบียน โดยมีจำนวนตัวอย่างในแต่ละคลาสตามตารางที่ \ref{tab:class-split}

\begin{table}[ht]
    \caption{จำนวนตัวอย่างในชุดฝึกและชุดทดสอบจำแนกตามคลาส}
    \label{tab:class-split}
    \small
    \begin{tabularx}{\linewidth}{@{}X r r r@{}}
        \hline
        \textbf{คลาส} & \textbf{ชุดฝึก} & \textbf{ชุดทดสอบ} & \textbf{รวม} \\
        \hline
        BENIGN & 1,450,574 & 621,680 & 2,072,254 \\
        DoS & 135,611 & 58,119 & 193,730 \\
        DDoS & 89,609 & 38,405 & 128,014 \\
        PortScan & 63,485 & 27,209 & 90,694 \\
        Brute Force & 6,405 & 2,745 & 9,150 \\
        Web Attack & 1,500 & 643 & 2,143 \\
        Bot & 1,363 & 585 & 1,948 \\
        Infiltration & 25 & 11 & 36 \\
        Heartbleed & 7 & 4 & 11 \\
        \hline
        \textbf{รวม} & \textbf{1,748,579} & \textbf{749,401} & \textbf{2,497,980} \\
        \hline
    \end{tabularx}
\end{table}

จากตารางที่ \ref{tab:class-split} เห็นได้ว่าคลาส BENIGN คิดเป็นสัดส่วนประมาณ 83.0\% ของชุดทดสอบ ซึ่งหมายความว่าแบบจำลองที่ทำนายว่าเป็น BENIGN ทุกกรณีจะได้ค่าความแม่นยำโดยรวมสูงถึง 0.8296 โดยไม่ตรวจจับการโจมตีได้เลยแม้แต่รายการเดียว ค่าดังกล่าวจึงถูกใช้เป็นแบบจำลองฐานเชิงเปรียบเทียบ (majority-class baseline) ในการรายงานผลทุกครั้ง นอกจากนี้ คลาส Heartbleed และ Infiltration มีจำนวนตัวอย่างในชุดทดสอบเพียงหลักหน่วยและหลักสิบตามลำดับ ค่าตัวชี้วัดรายคลาสของทั้งสองคลาสจึงต้องตีความอย่างระมัดระวังในฐานะค่าประมาณขอบเขตบน มิใช่ค่าประมาณที่เสถียร

\section{การพัฒนาแบบจำลอง}

\subsection{การกำหนดปัญหา}

ปัญหาถูกกำหนดให้เป็นการจำแนกประเภทแบบหลายคลาสที่มีจำนวนคลาส $K = 9$ โดยแบบจำลองรับเวกเตอร์คุณลักษณะ $\mathbf{x} \in \mathbb{R}^{77}$ และให้ผลลัพธ์เป็นเวกเตอร์ความน่าจะเป็นเหนือคลาสทั้งเก้า คลาสที่ถูกทำนายคือคลาสที่มีค่าความน่าจะเป็นสูงสุด ตามสมการ

\begin{equation}
    \hat{y} = \arg\max_{k \in \{1, \ldots, K\}} P(y = k \mid \mathbf{x})
\end{equation}

\subsection{แบบจำลองที่นำมาเปรียบเทียบ}

แบบจำลองที่นำมาเปรียบเทียบมี 7 รูปแบบ ครอบคลุมทั้งแบบจำลองเชิงเส้น แบบจำลองรวมกลุ่มชนิด Bagging และ Boosting โครงข่ายประสาทเทียม และการรวมกลุ่มเชิงซ้อน ตามที่สรุปไว้ในตารางที่ \ref{tab:model-summary} สำหรับ Stacking Ensemble โครงสร้างที่ใช้ประกอบด้วยแบบจำลองฐานสามตัว ได้แก่ LightGBM, XGBoost และ Random Forest โดยใช้ค่าความน่าจะเป็นของคลาสจากการทำนายนอกโฟลด์ (out-of-fold prediction) ด้วยการตรวจสอบไขว้ 5 โฟลด์ เป็นคุณลักษณะระดับเมทา และใช้ Logistic Regression เป็นแบบจำลองระดับที่สอง \cite{wolpert1992stacked, ting1999issues}

\subsection{การจัดการความไม่สมดุลของคลาส}

โครงงานนี้ใช้การถ่วงน้ำหนักคลาสแบบผกผันตามความถี่เป็นกลยุทธ์หลักและใช้แนวคิดเดียวกันในทุกแบบจำลอง เพื่อให้การจัดการความไม่สมดุลเป็นตัวแปรควบคุม โดยแบบจำลองส่วนใหญ่ (Logistic Regression, LightGBM, XGBoost และ Multi-Layer Perceptron) คำนวณน้ำหนักของคลาส $k$ จากสัดส่วนคลาสทั้งชุดฝึกครั้งเดียวตามสูตรมาตรฐานของ scikit-learn

\begin{equation}
    w_{k} = \frac{n}{K \cdot n_{k}}
\end{equation}

โดยที่ $n$ คือจำนวนตัวอย่างทั้งหมดในชุดฝึก $K$ คือจำนวนคลาส และ $n_{k}$ คือจำนวนตัวอย่างของคลาส $k$ ในชุดฝึก น้ำหนักนี้ถูกนำไปคูณกับพจน์ความสูญเสียของตัวอย่างในคลาสนั้น ๆ ทำให้ความผิดพลาดในคลาสส่วนน้อยมีต้นทุนสูงขึ้นตามสัดส่วน

โมเดลอีกสองรูปแบบใช้สูตรที่แปรผันตรงกับ $w_k$ ข้างต้นแต่คำนวณด้วยกลไกเฉพาะของแต่ละไลบรารี ได้แก่ Random Forest ใช้โหมด balanced\_subsample ซึ่งคำนวณน้ำหนักคลาสใหม่จากสัดส่วนคลาสภายในแต่ละตัวอย่างบูตสแตรป (bootstrap sample) ของต้นไม้แต่ละต้น แทนที่จะคำนวณครั้งเดียวจากทั้งชุดฝึก และ CatBoost ใช้โหมด auto\_class\_weights="Balanced" ซึ่งคำนวณจาก $w_k = \max_j(n_j) / n_k$ อันเป็นสูตรที่แปรผันตรงกับ $w_k$ ข้างต้นแต่ต่างค่าคงที่การนอร์มัลไลซ์ ทั้งสองกรณียังคงเป็นการถ่วงน้ำหนักแบบผกผันตามความถี่ในเชิงสัดส่วนระหว่างคลาสเหมือนกัน จึงไม่กระทบต่อการเปรียบเทียบระหว่างแบบจำลอง แต่ค่าคงที่ในการคำนวณมีความแตกต่างไปจากสมการข้างต้นตามข้อจำกัดของ API ในแต่ละไลบรารี

ประเด็นที่ต้องเน้นย้ำคือ การถ่วงน้ำหนักและการปรับสัดส่วนตัวอย่างใด ๆ กระทำเฉพาะภายในชุดฝึกและภายในโฟลด์ของการตรวจสอบไขว้เท่านั้น การกระจายตัวของคลาสในชุดทดสอบไม่ถูกแก้ไข เพื่อให้ผลการประเมินสะท้อนอัตราพื้นฐานที่เป็นจริง ซึ่งเป็นเงื่อนไขจำเป็นต่อการตีความอัตราการแจ้งเตือนผิดพลาดอย่างถูกต้อง \cite{axelsson2000baserate}

\subsection{การปรับไฮเปอร์พารามิเตอร์ภายใต้งบประมาณเดียวกัน}

เพื่อให้การเปรียบเทียบระหว่างแบบจำลองสะท้อนความสามารถของตัวแบบจำลองมากกว่าความพยายามในการปรับแต่ง โครงงานนี้กำหนดเงื่อนไขการค้นหาไฮเปอร์พารามิเตอร์ให้เหมือนกันใน 6 แบบจำลองที่ผ่านการค้นหา ได้แก่ Logistic Regression, Random Forest, XGBoost, LightGBM, CatBoost และ Multi-Layer Perceptron ดังนี้

\begin{itemize}
    \item วิธีการค้นหา: การค้นหาแบบสุ่ม (Randomized Search) ร่วมกับการตรวจสอบไขว้แบบแบ่งชั้น 5 โฟลด์
    \item ฟังก์ชันวัตถุประสงค์: ค่า Macro F1-score
    \item งบประมาณการค้นหา: 20 รอบการทดลองต่อแบบจำลอง บนตัวอย่างย่อยของชุดฝึกจำนวน 200,000 ระเบียน
    \item ขนาดปริภูมิการค้นหา: 144 ชุดค่าสำหรับทั้ง 6 แบบจำลอง ซึ่งทำให้สัดส่วนของปริภูมิที่ถูกสำรวจเท่ากันที่ประมาณ 13.9\%
\end{itemize}

การกำหนดขนาดปริภูมิการค้นหาให้เท่ากันมีความสำคัญ เพราะการใช้จำนวนรอบการทดลองเท่ากันบนปริภูมิที่มีขนาดต่างกันย่อมทำให้แบบจำลองที่มีปริภูมิเล็กกว่าได้รับการสำรวจในสัดส่วนที่สูงกว่า อย่างไรก็ตาม เงื่อนไขที่เท่ากันนี้ครอบคลุมเพียง 6 แบบจำลองข้างต้น Stacking Ensemble ไม่ได้ผ่านการค้นหาไฮเปอร์พารามิเตอร์ ทั้งในระดับแบบจำลองฐานและแบบจำลองระดับที่สอง เนื่องจากการค้นหาแต่ละรอบต้องฝึกแบบจำลองฐานทั้งชุดใหม่ทั้งหมด ทำให้ต้นทุนสูงเกินกว่าทรัพยากรที่มี แบบจำลองฐานของ Stacking จึงใช้ค่าตั้งต้นที่กำหนดไว้ในโค้ด มิใช่ค่าที่ได้จากการค้นหาในตารางที่ \ref{tab:best-params} ผลของ Stacking Ensemble จึงไม่ได้อยู่ภายใต้เงื่อนไขการปรับแต่งเดียวกับอีก 6 แบบจำลอง และควรตีความแยกจากการเปรียบเทียบภายใต้งบประมาณเท่ากัน ค่าไฮเปอร์พารามิเตอร์ที่ดีที่สุดของแต่ละแบบจำลองแสดงในตารางที่ \ref{tab:best-params}

\begin{table}[ht]
    \caption{ค่าไฮเปอร์พารามิเตอร์ที่ดีที่สุดของแต่ละแบบจำลองจากการค้นหาแบบสุ่ม}
    \label{tab:best-params}
    \small
    \begin{tabularx}{\linewidth}{@{}l >{\raggedright\arraybackslash}X@{}}
        \hline
        \textbf{แบบจำลอง} & \textbf{ค่าไฮเปอร์พารามิเตอร์ที่ดีที่สุด} \\
        \hline
        Logistic Regression & {C}~=~300.0, \texttt{tol}~=~$1\times10^{-4}$, \texttt{fit\_intercept}~=~True \\
        Random Forest & \texttt{n\_estimators}~=~300, \texttt{max\_depth}~=~30, \texttt{min\_samples\_split}~=~2, \texttt{min\_samples\_leaf}~=~1 \\
        XGBoost & \texttt{n\_estimators}~=~800, \texttt{max\_depth}~=~10, \texttt{learning\_rate}~=~0.05, \texttt{subsample}~=~0.8 \\
        LightGBM & \texttt{n\_estimators}~=~1000, \texttt{num\_leaves}~=~127, \texttt{learning\_rate}~=~0.1, \texttt{min\_child\_samples}~=~1 \\
        CatBoost & \texttt{iterations}~=~500, \texttt{depth}~=~6, \texttt{learning\_rate}~=~0.15, \texttt{l2\_leaf\_reg}~=~1 \\
        Multi-Layer Perceptron & \texttt{hidden\_layer\_sizes}~=~(256,~128,~64,~32), \texttt{learning\_rate\_init}~=~$5\times10^{-4}$, \texttt{batch\_size}~=~256, \texttt{alpha}~=~$1\times10^{-4}$ \\
        Stacking Ensemble & ไม่ได้ปรับไฮเปอร์พารามิเตอร์ (ใช้ค่าตั้งต้นในโค้ด) \\
        \hline
    \end{tabularx}
\end{table}

\section{การประเมินผลและเปรียบเทียบ}

\subsection{ตัวชี้วัดที่ใช้}

กำหนดให้ $TP_{k}$, $FP_{k}$, $TN_{k}$ และ $FN_{k}$ เป็นจำนวนตัวอย่างที่ทำนายถูกในคลาสบวก ทำนายผิดเป็นคลาสบวก ทำนายถูกในคลาสลบ และทำนายผิดเป็นคลาสลบ ของคลาส $k$ ตามลำดับ ค่าความเที่ยง ค่าความไว และค่า F1-score ของคลาส $k$ คำนวณได้จาก

\begin{equation}
    \mathrm{Precision}_{k} = \frac{TP_{k}}{TP_{k} + FP_{k}}, \qquad
    \mathrm{Recall}_{k} = \frac{TP_{k}}{TP_{k} + FN_{k}}
\end{equation}

\begin{equation}
    \mathrm{F1}_{k} = 2 \cdot \frac{\mathrm{Precision}_{k} \cdot \mathrm{Recall}_{k}}{\mathrm{Precision}_{k} + \mathrm{Recall}_{k}}
\end{equation}

ค่า Macro F1-score ซึ่งเป็นตัวชี้วัดหลักของโครงงานนี้ คือค่าเฉลี่ยเลขคณิตของ F1-score รายคลาสโดยให้น้ำหนักเท่ากันทุกคลาส ตามสมการ

\begin{equation}
    \mathrm{F1_{macro}} = \frac{1}{K} \sum_{k=1}^{K} \mathrm{F1}_{k}
\end{equation}

ในการทดลองหลักกำหนด $K=9$ และใช้ Macro F1 ของทั้ง 9 คลาสเป็นเกณฑ์จัดอันดับหลัก นอกจากนี้รายงาน Macro F1 เฉพาะคลาสที่มีตัวอย่างทดสอบไม่น้อยกว่า 30 ระเบียนเป็นตัวชี้วัดประกอบ ซึ่งในชุดทดสอบนี้เหลือ 7 คลาสหลังไม่นับ Heartbleed และ Infiltration ในค่าเฉลี่ย ส่วนการทดลองเสริมใช้ $K=7$ ตามโครงร่างคลาสร่วมของสองชุดข้อมูล ซึ่งรวม Infiltration แต่ไม่มี PortScan และ Heartbleed ดังนั้นตัวชี้วัด 7 คลาสทั้งสองกรณีจึงใช้คนละชุดคลาส

อัตราการแจ้งเตือนผิดพลาดและอัตราการพลาดตรวจจับคำนวณได้จาก

\begin{equation}
    \mathrm{FPR}_{k} = \frac{FP_{k}}{FP_{k} + TN_{k}}, \qquad
    \mathrm{FNR}_{k} = \frac{FN_{k}}{FN_{k} + TP_{k}}
\end{equation}

และค่าสัมประสิทธิ์สหสัมพันธ์ของแมทธิวส์สำหรับกรณีทวิภาค คำนวณได้จาก

\begin{equation}
    \mathrm{MCC} = \frac{TP \cdot TN - FP \cdot FN}{\sqrt{(TP+FP)(TP+FN)(TN+FP)(TN+FN)}}
\end{equation}

ซึ่งมีค่าอยู่ในช่วง $-1$ ถึง $1$ และทนทานต่อความไม่สมดุลของคลาสมากกว่าค่าความแม่นยำโดยรวม \cite{matthews1975comparison, chicco2020mcc}

\subsection{การประเมินผลสี่มิติ}

การประเมินแบบจำลองดำเนินการใน 4 มิติ เพื่อให้ครอบคลุมประเด็นที่เกี่ยวข้องกับการนำไปใช้งาน ดังนี้

\begin{enumerate}
    \item \textbf{ประสิทธิภาพการจำแนกประเภท (Classification Performance)} วัดด้วยค่า Accuracy, Balanced Accuracy, Precision, Recall และ F1-score ทั้งแบบ macro และ weighted รวมถึงค่า MCC โดยใช้ Macro F1-score เป็นเกณฑ์ตัดสินหลัก
    \item \textbf{ความสามารถในการตรวจจับการโจมตี (Attack Detection Ability)} ยุบผลการทำนายเป็นกรณีทวิภาคระหว่างการจราจรปกติและการโจมตี แล้ววัดอัตราการตรวจจับ ค่า FPR ค่า FNR และค่าความไวต่อการตรวจจับรายคลาส
    \item \textbf{ผลกระทบเชิงปฏิบัติการ (Operational Impact)} แปลงอัตราความผิดพลาดให้เป็นจำนวนสัมบูรณ์ ได้แก่ จำนวนการแจ้งเตือนผิดพลาดต่อทราฟฟิกปกติหนึ่งหมื่นระเบียน และจำนวนโฟลว์การโจมตีที่พลาดการตรวจจับ เพื่อสะท้อนภาระที่อาจเกิดขึ้นกับผู้ปฏิบัติงาน
    \item \textbf{ประสิทธิภาพเชิงคำนวณ (Computational Efficiency)} วัดเวลาที่ใช้ในการฝึก เวลาแฝงของการทำนายที่ควอนไทล์ที่ 50 และ 95 อัตราการประมวลผล และขนาดของไฟล์แบบจำลอง โดยการวัดเวลาแฝงดำเนินการบนหน่วยประมวลผลกลางสำหรับทุกแบบจำลองเป็นชุดข้อมูลขนาด 1,000 โฟลว์ ทำซ้ำ 30 ครั้งจากตำแหน่งเริ่มต้นแบบสุ่ม และตัดผลของรอบอุ่นเครื่องรอบแรกออก
\end{enumerate}

การวัดค่าเวลาแฝงที่ควอนไทล์ที่ 95 มีความสำคัญเทียบเท่าค่ามัชฌิม เนื่องจากระบบตรวจจับที่ทำงานเร็วเป็นส่วนใหญ่แต่ช้าเป็นครั้งคราวยังคงมีโอกาสสูญเสียโฟลว์ที่ประมวลผลไม่ทัน ซึ่งในทางปฏิบัติเทียบเท่ากับการพลาดตรวจจับ

\subsection{เกณฑ์การจัดอันดับแบบจำลอง}

เนื่องจากไม่มีตัวชี้วัดใดที่สามารถตัดสินความเหมาะสมของแบบจำลองได้ครบทุกมิติ โครงงานนี้จึงประกาศเกณฑ์การจัดอันดับไว้สามชุดควบคู่กัน แทนการเลือกแบบจำลองที่ดีที่สุดเพียงตัวเดียว ดังนี้

\begin{enumerate}
    \item \textbf{อันดับภาพรวม (Overall Best)} ตัดสินด้วยค่า Macro F1-score สูงสุด และใช้ค่า Weighted F1-score เป็นเกณฑ์ตัดสินกรณีเสมอ
    \item \textbf{อันดับด้านความมั่นคงปลอดภัย (Security-focused Best)} ตัดสินด้วยค่าความไวต่อการตรวจจับเฉลี่ยแบบ macro สูงสุด เฉพาะในกลุ่มแบบจำลองที่มีค่า FPR แบบทวิภาคไม่เกิน 0.01 และไม่มีคลาสใดที่มีค่าความไวเป็นศูนย์
    \item \textbf{อันดับด้านการนำไปใช้งาน (Deployment Best)} ตัดสินด้วยค่าเวลาแฝงที่ควอนไทล์ที่ 95 ต่ำสุด เฉพาะในกลุ่มแบบจำลองที่มีค่า Macro F1-score ต่างจากแบบจำลองที่ดีที่สุดไม่เกิน 0.02 และใช้ขนาดของไฟล์แบบจำลองเป็นเกณฑ์ตัดสินกรณีเสมอ
\end{enumerate}

เกณฑ์ทั้งสามชุดถูกระบุไว้ก่อนการวิเคราะห์ผล (criteria specified before analysis) ในไฟล์นโยบายการจัดอันดับ \texttt{configs/ranking\_policy.json} และนำไปใช้กับทุกแบบจำลองด้วยกฎเดียวกัน เพื่อลดโอกาสการเลือกเกณฑ์ให้เข้ากับผลลัพธ์ภายหลัง

\subsection{การตรวจสอบความน่าเชื่อถือของผลลัพธ์}

เพื่อตรวจว่าค่าตัวชี้วัดที่สูงอาจเกิดจากการรั่วไหลของข้อมูลหรือความบังเอิญของการแบ่งข้อมูลเพียงครั้งเดียว โครงงานนี้กำหนดการตรวจสอบเพิ่มเติมสามรายการ ประกอบด้วย การเปรียบเทียบกับแบบจำลองฐานที่ทำนายคลาสส่วนใหญ่เสมอ การรายงานค่าเฉลี่ยและส่วนเบี่ยงเบนมาตรฐานของ Macro F1-score จากการตรวจสอบไขว้ 5 โฟลด์บนชุดฝึก และการรายงานค่า Macro F1-score เฉพาะกลุ่มคลาสที่มีจำนวนตัวอย่างในชุดทดสอบไม่น้อยกว่า 30 ระเบียน ควบคู่กับค่า Macro F1-score ของทุกคลาส เพื่อให้ผู้อ่านแยกแยะได้ว่าค่าตัวชี้วัดส่วนใดได้รับผลกระทบจากคลาสที่มีตัวอย่างน้อยเกินกว่าจะประมาณค่าได้อย่างเสถียร

\section{การอธิบายผลของแบบจำลอง}
\label{sec:shap}

การวิเคราะห์เพื่ออธิบายผลของแบบจำลองดำเนินการด้วยเทคนิค SHAP ตามหลักการที่กล่าวถึงในหัวข้อ \ref{sec:xai-review} โดยใช้อัลกอริทึม TreeSHAP สำหรับแบบจำลองที่อาศัยต้นไม้ตัดสินใจ ได้แก่ Random Forest, XGBoost, LightGBM และ CatBoost ซึ่งให้ผลลัพธ์ที่แม่นยำในเวลาแบบพหุนาม \cite{lundberg2020treeshap} และใช้ Kernel SHAP ซึ่งเป็นวิธีประมาณค่าแบบไม่ขึ้นกับชนิดของแบบจำลองสำหรับ Multi-Layer Perceptron โดยคำนวณบนตัวอย่างย่อยเนื่องจากมีต้นทุนการคำนวณสูง

ผลการวิเคราะห์ถูกนำเสนอในสองระดับ ระดับแรกคือระดับภาพรวม (global explanation) ซึ่งจัดอันดับคุณลักษณะตามค่าเฉลี่ยของค่าสัมบูรณ์ของค่า SHAP เพื่อระบุว่าคุณลักษณะใดมีอิทธิพลต่อการตัดสินใจของแบบจำลองมากที่สุดโดยรวม และระดับที่สองคือระดับรายตัวอย่าง (local explanation) ซึ่งแสดงว่าคุณลักษณะแต่ละตัวผลักผลการทำนายของโฟลว์หนึ่งระเบียนไปในทิศทางใดและด้วยขนาดเท่าใด นอกจากการอธิบายผลแล้ว การวิเคราะห์นี้ยังทำหน้าที่ตรวจสอบด้วยว่าแบบจำลองเรียนรู้จากคุณลักษณะที่สอดคล้องกับลักษณะทางเทคนิคของการโจมตีจริง มากกว่าจากลักษณะเฉพาะของกระบวนการสร้างชุดข้อมูล

\section{การแสดงผลและระบบต้นแบบ}

ผลลัพธ์ทั้งหมดถูกนำเสนอผ่านช่องทางต่อไปนี้

\begin{enumerate}
    \item ตารางเปรียบเทียบตัวชี้วัดของแบบจำลองทั้ง 7 รูปแบบในทั้งสี่มิติของการประเมิน
    \item เมทริกซ์ความสับสน (Confusion Matrix) แบบจำนวนนับของแบบจำลองทั้ง 7 รูปแบบจากการทดลองหลักบน CICIDS2017 โดยแสดง CatBoost ในหัวข้อ \ref{sec:main-classification} และอีก 6 แบบจำลองในภาคผนวก \ref{appendix:main-confusion-matrices}
    \item แผนภาพสรุปผลการวิเคราะห์ด้วย SHAP ทั้งในระดับภาพรวมและระดับรายตัวอย่าง
    \item ระบบต้นแบบในรูปเว็บแดชบอร์ด พัฒนาด้วยเฟรมเวิร์ก FastAPI สำหรับส่วนบริการเบื้องหลัง และ Next.js สำหรับส่วนติดต่อผู้ใช้ ซึ่งรองรับการเรียกดูผลการเปรียบเทียบแบบจำลอง การตรวจสอบการกระจายตัวของข้อมูล การเรียกดูผลการอธิบายแบบจำลอง และการอัปโหลดไฟล์ข้อมูลโฟลว์ใหม่เพื่อทำนายผลเป็นชุด
\end{enumerate}

ส่วนของการทำนายผลข้อมูลใหม่มีการตรวจสอบโครงสร้างของไฟล์ข้อมูลนำเข้า (schema validation) โดยเปรียบเทียบกับรายชื่อคุณลักษณะที่บันทึกไว้พร้อมกับไฟล์แบบจำลอง หากไฟล์ข้อมูลนำเข้าขาดคุณลักษณะที่จำเป็น ระบบจะปฏิเสธการทำนายและแจ้งรายชื่อคุณลักษณะที่ขาดไป เพื่อป้องกันการให้ผลลัพธ์ที่ไม่ถูกต้องจากข้อมูลนำเข้าที่ไม่สอดคล้องกับที่แบบจำลองได้รับการฝึกไว้

\section{การทดลองเสริม: การส่งผ่านข้ามชุดข้อมูลแบบสองทิศทาง}
\label{sec:crossdataset}

\subsection{คำถามของการทดลองและเหตุผลของการออกแบบ}

การประเมินผลทั้งหมดในหัวข้อก่อนหน้ากระทำอยู่ภายในชุดข้อมูล CICIDS2017 ซึ่งบันทึกจากสภาพแวดล้อมเครือข่ายจำลองเพียงชุดเดียว แม้การแบ่งข้อมูลตามลำดับการจับข้อมูลจะช่วยลดความเสี่ยงของการรั่วไหลของข้อมูลภายในคอร์ปัส แต่ยังไม่สามารถตอบได้ว่าแบบจำลองเรียนรู้พฤติกรรมของการโจมตีที่ใช้ได้ทั่วไป หรือเรียนรู้ลักษณะเฉพาะของสภาพแวดล้อมการจับข้อมูลนั้น โครงงานนี้จึงจัดการทดลองเสริมขึ้นเพื่อตอบคำถามว่า \emph{แบบจำลองที่ฝึกจากคอร์ปัสหนึ่งยังคงใช้งานได้ดีเพียงใดเมื่อนำไปทำนายข้อมูลจากอีกคอร์ปัสหนึ่งที่บันทึกจากสภาพแวดล้อมต่างกัน} โดยใช้ชุดข้อมูล CSE-CIC-IDS2018 เป็นคอร์ปัสคู่เปรียบเทียบ

ค่า Macro F1-score ที่วัดได้จากคอร์ปัสอื่นเพียงค่าเดียวไม่สามารถตีความได้โดยลำพัง เพราะไม่ทราบว่าค่าที่ต่ำนั้นเกิดจากการส่งผ่านที่ล้มเหลว หรือเกิดจากงานจำแนกประเภทบนคอร์ปัสนั้นยากกว่าโดยตัวมันเอง การตีความจึงต้องอาศัยค่าอ้างอิงสองค่าบนชุดทดสอบเดียวกัน ได้แก่ \textbf{ค่าพื้น} (floor) ซึ่งเป็นค่าของแบบจำลองฐานที่ทำนายคลาสส่วนใหญ่เสมอ และ \textbf{ค่าเพดาน} (ceiling) ซึ่งเป็นค่าที่แบบจำลองชนิดเดียวกันทำได้เมื่อฝึกและทดสอบภายในคอร์ปัสปลายทางเอง ด้วยเหตุนี้การทดลองจึงต้องฝึกแบบจำลองจากทั้งสองคอร์ปัส มิใช่ใช้ CSE-CIC-IDS2018 เป็นเพียงชุดทดสอบภายนอก ผังการทดลองจึงมีลักษณะเป็นเมทริกซ์ 2×2 ตามตารางที่ \ref{tab:crossdataset-design}

\begin{table}[ht]
    \caption{ผังการทดลองแบบสองทิศทางและบทบาทของผลการวัดในแต่ละช่อง}
    \label{tab:crossdataset-design}
    \small
    \begin{tabularx}{\linewidth}{@{}l l X@{}}
        \hline
        \textbf{คอร์ปัสที่ใช้ฝึก} & \textbf{คอร์ปัสที่ใช้ทดสอบ} & \textbf{บทบาทของผลการวัด} \\
        \hline
        CICIDS2017 & CICIDS2017 & ค่าเพดานของทิศทาง 2018\dirarrow{}2017 \\
        CICIDS2017 & CSE-CIC-IDS2018 & การส่งผ่านทิศทาง 2017\dirarrow{}2018 \\
        CSE-CIC-IDS2018 & CSE-CIC-IDS2018 & ค่าเพดานของทิศทาง 2017\dirarrow{}2018 \\
        CSE-CIC-IDS2018 & CICIDS2017 & การส่งผ่านทิศทาง 2018\dirarrow{}2017 \\
        \hline
    \end{tabularx}
\end{table}

แบบจำลองในการทดลองเสริมนี้ถูกฝึกขึ้นใหม่แยกจากแบบจำลองหลักของโครงงาน โดยใช้โครงร่างคลาสร่วม 7 คลาส ตัวอย่างย่อย และค่าไฮเปอร์พารามิเตอร์ที่ไม่ผ่านการปรับแต่ง ค่าตัวชี้วัดในหัวข้อนี้จึงไม่สามารถนำไปเทียบตรงกับผลการทดลองหลักแบบ 9 คลาสได้ และไม่มีผลต่อการคัดเลือกแบบจำลองแชมเปียนที่รายงานในบทที่ \ref{chapter:result}

\subsection{ชุดข้อมูลและการปรับให้เทียบกันได้}

ชุดข้อมูล CICIDS2017 หลังการทำความสะอาดมี 2,497,980 ระเบียน (หัวข้อ \ref{sec:split}) ส่วนชุดข้อมูล CSE-CIC-IDS2018 \cite{cic2018web} ประกอบด้วยไฟล์บันทึกการจับข้อมูลรายวัน 10 ไฟล์ รวม 16,232,943 ระเบียน และเหลือ 11,445,955 ระเบียนหลังการทำความสะอาดด้วยกฎชุดเดียวกัน เนื่องจากทั้งสองคอร์ปัสมีชื่อคอลัมน์คุณลักษณะและชุดป้ายกำกับไม่ตรงกัน จึงต้องปรับให้อยู่ในรูปแบบร่วมกัน (canonicalization) ตามลำดับดังนี้

\begin{enumerate}
    \item \textbf{การจับคู่คุณลักษณะ} คุณลักษณะร่วม 77 ตัวถูกจับคู่ตามความหมายของค่าที่ CICFlowMeter สร้าง มิใช่การเทียบชื่อคอลัมน์ เนื่องจาก 50 จาก 77 คอลัมน์มีชื่อต่างกันระหว่างสองคอร์ปัส ความถูกต้องของการจับคู่ได้รับการตรวจสอบด้วยความสัมพันธ์เชิงพีชคณิตระหว่างคอลัมน์ สเกลของค่า และชนิดข้อมูล ไม่พบคู่ใดที่บ่งชี้ว่าจับคู่ผิด แม้บางคู่จะมีสเกลค่าต่างกันมาก (เช่น Fwd IAT Mean ต่างกันประมาณ 4.5 อันดับของเลขยกกำลังสิบ)
    \item \textbf{การรวมคลาส} ป้ายกำกับของทั้งสองคอร์ปัสถูกรวบเข้าเป็นคลาสร่วม 7 คลาส ได้แก่ Benign, DDoS, DoS, Brute Force, Bot, Infiltration และ Web Attack คลาส PortScan และ Heartbleed ของ CICIDS2017 ถูกตัดออกเพราะไม่มีคลาสคู่กันใน CSE-CIC-IDS2018 ทำให้ CICIDS2017 เหลือ 2,407,275 ระเบียน
    \item \textbf{ลำดับขั้นตอน} ทั้งสองคอร์ปัสผ่านขั้นตอนเดียวกันตามลำดับ คือ เปลี่ยนชื่อคอลัมน์เป็นรูปแบบร่วม รวมคลาสและตัดแถวที่ไม่มีคลาสคู่กัน ทำความสะอาด แล้วจึงสุ่มตัวอย่าง การทำความสะอาดก่อนการสุ่มทำให้ตัวอย่างทุกแถวที่สุ่มได้เป็นแถวที่ใช้งานได้จริงทั้งสองฝั่ง
    \item \textbf{การสุ่มตัวอย่าง} สุ่มตัวอย่างคอร์ปัสละ 300,000 ระเบียนด้วยค่าเริ่มต้นการสุ่ม 42 โดยเก็บทุกแถวของคลาสที่มีน้อยกว่า 10,000 ระเบียน แล้วแบ่งจำนวนที่เหลือให้คลาสใหญ่ตามสัดส่วนเดิมของคอร์ปัส วิธีนี้รักษาทั้งคลาสหายากและสัดส่วนความไม่สมดุลเดิมไว้ (Benign ประมาณ 83--88\% ทั้งสองฝั่ง) ตัวอย่างชุดนี้สุ่มครั้งเดียวและใช้ร่วมกันในทุกรอบของการทดลอง
\end{enumerate}

ตัวอย่างฝั่ง CICIDS2017 ประกอบด้วย Benign 248,188 ระเบียน DoS 23,203 ระเบียน DDoS 15,332 ระเบียน Brute Force 9,150 ระเบียน Web Attack 2,143 ระเบียน Bot 1,948 ระเบียน และ Infiltration 36 ระเบียน ขณะที่ฝั่ง CSE-CIC-IDS2018 ประกอบด้วย Benign 263,858 ระเบียน DDoS 20,272 ระเบียน DoS 5,132 ระเบียน Bot 3,778 ระเบียน Infiltration 3,642 ระเบียน Brute Force 2,459 ระเบียน และ Web Attack 859 ระเบียน คลาส Infiltration ซึ่งมีตัวอย่างเพียง 36 ระเบียนใน CICIDS2017 กลับมีตัวอย่างหลักพันใน CSE-CIC-IDS2018 ความแตกต่างของความถี่รายคลาสนี้เป็นปัจจัยที่ต้องนำมาประกอบการตีความผล

\subsection{การแบ่งชุดฝึกและชุดทดสอบ}
\label{sec:cd-split}

โพรโทคอลสุดท้ายของการทดลองเสริม (เรียกว่า protocol\_v3) แบ่งตัวอย่าง 300,000 ระเบียนของแต่ละคอร์ปัสเป็นชุดฝึก 70\% และชุดทดสอบ 30\% โดยดำเนินการสองรูปแบบคู่ขนานกัน

\begin{enumerate}
    \item \textbf{การแบ่งแบบสุ่ม (random split)} แบ่งแบบแบ่งชั้นตามคลาส (stratified) โดยสลับลำดับแถวภายในแต่ละคลาสด้วยค่าเริ่มต้นการสุ่มของรอบนั้น
    \item \textbf{การแบ่งตามลำดับการจับข้อมูล (capture-order split)} ภายในแต่ละกลุ่ม (กลุ่มการจับข้อมูล, คลาส) เรียงแถวตามกุญแจลำดับ แล้วจัด 70\% แรกเป็นชุดฝึกและ 30\% ท้ายเป็นชุดทดสอบ สำหรับ CICIDS2017 กลุ่มการจับข้อมูลคือไฟล์ต้นฉบับ (source\_file) และกุญแจลำดับคือตำแหน่งของแถวในไฟล์ CSV ต้นฉบับ (\_row\_index) เนื่องจากชุดข้อมูลรุ่น \textit{MachineLearningCVE} ไม่มีคอลัมน์เวลา ส่วน CSE-CIC-IDS2018 กลุ่มการจับข้อมูลคือวันที่บันทึก และกุญแจลำดับคือคอลัมน์ Timestamp
\end{enumerate}

ต้องระบุให้ชัดว่า การแบ่งตามลำดับการจับข้อมูลของ CICIDS2017 \emph{ไม่ใช่การแบ่งตามเวลาประทับ (timestamp split)} แต่เป็นการแบ่งตามลำดับแถวในไฟล์ ซึ่งได้รับการตรวจสอบแล้วว่าสอดคล้องกับตารางเวลาการโจมตีที่ผู้จัดทำชุดข้อมูลเผยแพร่ และการแบ่งเกิดขึ้นภายในแต่ละกลุ่ม (ไฟล์, คลาส) มิใช่การตัดที่เวลาเดียวกันทั้งคอร์ปัส เพราะการโจมตีแต่ละคลาสปรากฏอยู่ในไฟล์เดียว การตัดที่เวลาเดียวทั้งคอร์ปัสจะทำให้บางคลาสไปอยู่ฝั่งเดียวทั้งหมด ผลที่ตามมาคือ ภายในคลาสที่ประกอบด้วยเครื่องมือโจมตีหลายตัวที่ทำงานต่อเนื่องกัน ตำแหน่งที่ตัดจะกำหนดว่าเครื่องมือใดอยู่ในชุดฝึก ซึ่งส่งผลต่อผลลัพธ์อย่างมาก (หัวข้อ \ref{sec:cd-ssh})

การแบ่งตามลำดับการจับข้อมูลเป็นฟังก์ชันของข้อมูลล้วน ๆ จึงได้ชุดฝึกและชุดทดสอบเดิมทุกรอบ ส่วนการแบ่งแบบสุ่มเปลี่ยนตามค่าเริ่มต้นการสุ่ม ในโพรโทคอลก่อนหน้า (protocol\_v2) สัดส่วนการแบ่งคือ 60/10/30 โดยส่วน 10\% ถูกกันไว้สำหรับการปรับเทียบแต่ไม่ได้ถูกใช้งาน โครงงานจึงตัดส่วนดังกล่าวออกและใช้ 70/30 เป็นโพรโทคอลสุดท้าย ผลของ protocol\_v2 ถูกอ้างถึงในบทที่ \ref{chapter:result} เฉพาะเพื่อวิเคราะห์ความไวต่อสัดส่วนการแบ่งเท่านั้น

\subsection{แบบจำลอง รอบการทดลอง และตัวชี้วัด}

ในแต่ละรอบของการทดลอง (กำหนดโดยรูปแบบการแบ่งข้อมูลและค่าเริ่มต้นของการสุ่ม) แบบจำลองแต่ละรูปแบบจะถูกฝึกจากชุดฝึกของ CICIDS2017 และจากชุดฝึกของ CSE-CIC-IDS2018 แยกกัน แล้วนำแต่ละชุดไปทดสอบกับชุดทดสอบของ \emph{ทั้งสอง} คอร์ปัส ตัวปรับสเกลคุณลักษณะถูกประมาณค่าจากคอร์ปัสต้นทางเพียงฝั่งเดียวแล้วนำไปใช้กับชุดทดสอบของอีกฝั่งโดยไม่ปรับค่าใหม่ ซึ่งตรงกับกรณีที่นำแบบจำลองที่ฝึกเสร็จแล้วไปใช้กับข้อมูลของอีกสภาพแวดล้อมหนึ่งโดยไม่มีข้อมูลของสภาพแวดล้อมนั้นมาปรับเทียบ เงื่อนไขของการทดลองมีดังนี้

\begin{itemize}
    \item แบบจำลอง 7 รูปแบบเดียวกับการทดลองหลัก ได้แก่ Logistic Regression, Random Forest, XGBoost, LightGBM, CatBoost, Multi-Layer Perceptron และ Stacking Ensemble ฝึกขึ้นใหม่จากศูนย์ด้วยค่าไฮเปอร์พารามิเตอร์ตั้งต้นชุดเดียวกันทั้งสองทิศ พร้อมการถ่วงน้ำหนักคลาสแบบผกผันตามความถี่ โดยจงใจไม่นำค่าที่ปรับแต่งไว้แล้วบน CICIDS2017 (ตารางที่ \ref{tab:best-params}) มาใช้ เพราะจะทำให้ฝั่ง CICIDS2017 ได้เปรียบ
    \item ค่าเริ่มต้นการสุ่ม 5 ค่า (42--46) กับการแบ่งข้อมูล 2 รูปแบบ ได้ผลทั้งสิ้น 280 ผลการวัด (2 รูปแบบ × 5 ค่าเริ่มต้น × 2 คอร์ปัสฝึก × 2 คอร์ปัสทดสอบ × 7 แบบจำลอง)
    \item ภายใต้การแบ่งตามลำดับการจับข้อมูล ชุดฝึกและชุดทดสอบเหมือนกันทุกรอบ ส่วนเบี่ยงเบนมาตรฐานระหว่าง 5 รอบจึงวัดเฉพาะความสุ่มภายในตัวแบบจำลอง มิใช่ความผันผวนของการแบ่งข้อมูล
\end{itemize}

ตัวชี้วัดหลักคือ Macro F1-score ของทั้ง 7 คลาส รายงานเป็นค่าเฉลี่ย ± ส่วนเบี่ยงเบนมาตรฐานจาก 5 รอบ ค่าพื้นของแบบจำลองที่ทำนาย Benign ทั้งหมดเท่ากับ 0.1294 บนชุดทดสอบของ CICIDS2017 และ 0.1337 บนชุดทดสอบของ CSE-CIC-IDS2018 เพื่อให้การส่งผ่านทั้งสองทิศเทียบกันได้ทั้งที่ค่าเพดานและค่าพื้นของแต่ละคอร์ปัสไม่เท่ากัน จึงนิยามสัดส่วนประสิทธิภาพที่กู้คืนได้ (Recovery) ดังสมการ

\begin{equation}
    \mathrm{Recovery} = \frac{\mathrm{F1_{transfer}} - \mathrm{F1_{floor}}}{\mathrm{F1_{ceiling}} - \mathrm{F1_{floor}}} \times 100\%
\end{equation}

โดยที่ F1\textsubscript{transfer} คือค่าที่แบบจำลองซึ่งฝึกจากคอร์ปัสต้นทางทำได้บนชุดทดสอบของคอร์ปัสปลายทาง F1\textsubscript{floor} คือค่าพื้นบนชุดทดสอบปลายทาง และ F1\textsubscript{ceiling} คือค่าที่แบบจำลองชนิดเดียวกันทำได้เมื่อฝึกและทดสอบภายในคอร์ปัสปลายทาง ภายใต้รูปแบบการแบ่งข้อมูลเดียวกัน ค่า 100\% หมายถึงการส่งผ่านที่ไม่สูญเสียประสิทธิภาพเลย และ 0\% หมายถึงดีไม่ต่างจากการเดาคลาสส่วนใหญ่ เนื่องจากค่าเพดานแตกต่างกันในแต่ละแบบจำลอง ค่า Recovery จึงต้องอ่านคู่กับค่า Macro F1-score ดิบเสมอ และไม่ใช้จัดอันดับแบบจำลองโดยลำพัง ความไม่สมมาตรของการส่งผ่าน (asymmetry) นิยามเป็นผลต่างของค่า Recovery เฉลี่ย 7 แบบจำลองระหว่างทิศ 2018\dirarrow{}2017 กับทิศ 2017\dirarrow{}2018 มีหน่วยเป็นจุดเปอร์เซ็นต์ นอกจากนี้ยังรายงานค่า F1 รายคลาสเพื่อระบุว่าคะแนนที่เปลี่ยนไปมาจากคลาสใด

\subsection{การตรวจสอบร่องรอยเฉพาะของชุดข้อมูล}
\label{sec:cd-av-method}

เพื่อตรวจว่าสองคอร์ปัสแยกออกจากกันได้ด้วยค่าของคุณลักษณะมากเพียงใด (dataset fingerprint) โครงงานใช้การตรวจสอบเชิงปฏิปักษ์ (adversarial validation) โดยฝึกตัวจำแนก LightGBM ให้ทายว่าระเบียนหนึ่งมาจากคอร์ปัสใด แยกทีละคลาสร่วม ใช้ตัวอย่างไม่เกิน 8,000 ระเบียนต่อฝั่ง ทำซ้ำ 5 ค่าเริ่มต้นการสุ่ม และวัดด้วยค่า ROC AUC ผลดังกล่าวมีการตรวจควบคุมประกอบดังนี้

\begin{itemize}
    \item \textbf{ลำดับความจุของแบบจำลอง (model ladder)} ทำซ้ำด้วย decision stump (กฎเดียวบนคอลัมน์เดียว) ต้นไม้ลึก 3 ชั้น และ Logistic Regression เพื่อดูว่าความสามารถในการแยกต้องอาศัยแบบจำลองที่ซับซ้อนหรือไม่
    \item \textbf{การสลับป้ายกำกับ (label-shuffle)} ใช้แถวจริงชุดเดียวกันแต่สลับป้ายกำกับคอร์ปัสก่อนการแบ่งข้อมูล หากยังแยกได้แสดงว่าไปป์ไลน์รับรู้ป้ายกำกับผ่านช่องทางอื่นที่ไม่ใช่คุณลักษณะ
    \item \textbf{กลุ่มควบคุมว่าง (null control)} แบ่งคอร์ปัสเดียวออกเป็นสองครึ่งแบบสุ่ม เพื่อตรวจว่าวิธีการสร้างการแยกขึ้นเองหรือไม่
    \item \textbf{กลุ่มควบคุมภายในคอร์ปัส} แยกคนละไฟล์หรือคนละช่วงเวลาภายในคอร์ปัสเดียวกัน ด้วยแบบจำลองทั้ง 7 รูปแบบและจำนวนแถวต่อฝั่งเท่ากัน และกรณีที่ตรึงเครื่องมือโจมตีให้เป็นตัวเดียวกันแล้วเปลี่ยนเฉพาะคอร์ปัส
    \item \textbf{ระเบียนซ้ำ} นับระเบียนที่ซ้ำกันทุกค่าคร่อมชุดฝึกและชุดทดสอบ และระเบียนใกล้ซ้ำ (near-duplicate) ซึ่งนิยามเป็นแถวทดสอบที่อยู่ใกล้แถวฝึกมากกว่าระยะที่แถวฝึกอยู่ห่างกันเองถึง 99.9\% ของกรณี ระดับที่เกิดขึ้นโดยบังเอิญจึงเท่ากับ 0.10\%
\end{itemize}

\subsection{การวิเคราะห์สาเหตุของการทำนายพลาดด้วย SHAP}
\label{sec:cd-shap-method}

สมมติฐานที่ทดสอบคือ \emph{เมื่อนำแบบจำลองไปใช้ข้ามคอร์ปัส ระเบียนโจมตีที่ถูกทำนายผิดเป็น Benign ถูกผลักไปทาง Benign ผ่านกลุ่มคุณลักษณะที่บ่งชี้คอร์ปัสมากกว่าระเบียนที่ทำนายถูก} กลุ่มคุณลักษณะดังกล่าว (F) เลือกจากคุณลักษณะ 60 ตัว (77 คอลัมน์หลังรวมคอลัมน์ที่มีค่าซ้ำกัน) โดยใช้ผลการตรวจสอบเชิงปฏิปักษ์ และนับเฉพาะคุณลักษณะที่ผ่านทั้งเกณฑ์ permutation importance และเกณฑ์ SHAP ได้กลุ่มหลัก (consensus) 13 ตัว ได้แก่ min\_seg\_size\_forward, RST Flag Count, Init\_Win\_bytes\_forward, Destination Port, Max Packet Length, act\_data\_pkt\_fwd, Fwd IAT Min, Init\_Win\_bytes\_backward, Flow IAT Min, Bwd IAT Min, Min Packet Length, Bwd Packet Length Min และ Bwd Packet Length Std กลุ่มนี้ถูกกำหนดไว้ก่อนการทดสอบและไม่ถูกเปลี่ยนภายหลัง

การวิเคราะห์ใช้แบบจำลองต้นไม้ 4 รูปแบบ (LightGBM, Random Forest, XGBoost และ CatBoost) บนการแบ่งตามลำดับการจับข้อมูล 70/30 ค่าเริ่มต้นการสุ่ม 42--46 และใช้ไม่เกิน 1,000 แถวต่อคลาสต่อกลุ่มต่อรอบ สำหรับแต่ละแถวคำนวณ \textit{s} คือสัดส่วนของค่า SHAP ที่เป็นบวกต่อผลลัพธ์ Benign ซึ่งมาจากคุณลักษณะใน F การเปรียบเทียบในเกณฑ์ตัดสินสุดท้ายเป็นแบบ \textbf{ภายในคอร์ปัสปลายทาง} (within-target) คือ

\begin{equation}
    G_{\mathrm{within}} = \bar{s}(\mathrm{missed}) - \bar{s}(\mathrm{caught_{tgt}})
\end{equation}

โดยที่ missed คือระเบียนโจมตีในชุดทดสอบของคอร์ปัสปลายทางที่ถูกทำนายเป็น Benign และ caught\textsubscript{tgt} คือระเบียนโจมตีคลาสเดียวกันในคอร์ปัสเดียวกันที่ถูกทำนายถูก การเทียบภายในคอร์ปัสปลายทางตัดผลของการที่ข้อมูลมาจากต่างคอร์ปัสออก ซึ่งการเทียบกับระเบียนที่ทำนายถูกในคอร์ปัสต้นทาง (ที่ใช้ในรอบก่อนหน้า) แยกไม่ได้ ค่า \textit{G} ถูกเทียบกับกลุ่มคุณลักษณะสุ่มขนาดเท่ากัน 2,000 กลุ่ม (\textit{p} แบบด้านเดียว) กรณีที่นับได้ต้องมีระเบียน missed อย่างน้อย 100 แถว และ caught\textsubscript{tgt} อย่างน้อย 30 แถวในทุกรอบ และกรณีหนึ่งจะถือว่าผ่านเมื่อ \textit{G}\textsubscript{within} > 0 และ \textit{p} ต่ำกว่า 0.05 หารด้วยจำนวนกรณีที่นับได้ของแบบจำลองนั้น (Bonferroni) เกณฑ์ทั้งหมดถูกระบุไว้ก่อนการวิเคราะห์ (criteria specified before analysis) ในเอกสารภายในของโครงงานที่ระบุวันที่ไว้ก่อนการรัน แต่มิได้ลงทะเบียนกับหน่วยงานหรือระบบภายนอก จึงไม่ถือเป็นการลงทะเบียนล่วงหน้า (pre-registration) ในความหมายที่เป็นทางการ

\section{การควบคุมให้การทดลองทำซ้ำได้}

เพื่อให้ผลการทดลองสามารถทำซ้ำและตรวจสอบย้อนกลับได้ โครงงานนี้กำหนดมาตรการควบคุมดังนี้

\begin{enumerate}
    \item สำหรับการทดลองหลัก กำหนดค่าเริ่มต้นของการสุ่ม (random seed) เท่ากับ 42 เพื่อให้สามารถทำซ้ำผลได้ ส่วนการทดลองเสริมข้ามชุดข้อมูลใช้ค่าเริ่มต้นการสุ่ม 5 ค่า ได้แก่ 42--46 เพื่อวัดความแปรปรวนจากการสุ่มของแบบจำลองและการแบ่งข้อมูลแบบสุ่ม ตามที่อธิบายในหัวข้อ 3.8.4 สำหรับการแบ่งตามลำดับการจับข้อมูล ชุดฝึกและชุดทดสอบคงเดิมทั้ง 5 รอบ จึงวัดความแปรปรวนจากความสุ่มของแบบจำลอง ส่วนการสุ่มตัวอย่างคอร์ปัสละ 300,000 ระเบียนใช้ seed 42 เพียงครั้งเดียวและใช้ตัวอย่างชุดเดิมทุกรอบ
    \item ตรึงรุ่นของไลบรารีที่ใช้ทั้งหมดไว้ในไฟล์รายการที่ต้องติดตั้ง
    \item คำนวณค่าลายนิ้วมือดิจิทัลแบบ SHA-256 ของชุดข้อมูลที่ผ่านการทำความสะอาดแล้ว และบันทึกไว้พร้อมกับผลลัพธ์ของทุกการทดลอง เพื่อยืนยันว่าผลลัพธ์แต่ละชุดมาจากข้อมูลชุดเดียวกัน
    \item บันทึกไฟล์ประกาศการแบ่งข้อมูล (split manifest) ซึ่งระบุจำนวนตัวอย่างที่คาดหมายในแต่ละคลาส และมีขั้นตอนตรวจสอบที่สร้างการแบ่งข้อมูลขึ้นใหม่จากแคชแล้วเปรียบเทียบกับค่าที่ประกาศไว้
    \item บันทึกรายชื่อและลำดับของคุณลักษณะไว้พร้อมกับไฟล์แบบจำลองแต่ละไฟล์ เพื่อป้องกันความคลาดเคลื่อนของโครงสร้างข้อมูลระหว่างขั้นตอนการฝึกและขั้นตอนการทำนาย
    \item จัดทำชุดทดสอบซอฟต์แวร์แบบอัตโนมัติครอบคลุมไปป์ไลน์การประมวลผลข้อมูล ส่วนบริการเบื้องหลัง และวงจรชีวิตของไฟล์ผลลัพธ์
\end{enumerate}
